%%%%%%%%%%%%%%%%%%%%%%%%%%%%%%%%%%%%%%%%%
% Medium Length Professional CV
% LaTeX Template
% Version 2.0 (8/5/13)
%
% This template has been downloaded from:
% http://www.LaTeXTemplates.com
%
% Original author:
% Trey Hunner (http://www.treyhunner.com/)
%
% Important note:
% This template requires the resume.cls file to be in the same directory as the
% .tex file. The resume.cls file provides the resume style used for structuring the
% document.
%
%%%%%%%%%%%%%%%%%%%%%%%%%%%%%%%%%%%%%%%%%

%----------------------------------------------------------------------------------------
%	PACKAGES AND OTHER DOCUMENT CONFIGURATIONS
%----------------------------------------------------------------------------------------

\documentclass{resume} % Use the custom resume.cls style

\usepackage[left=0.6in,top=0.5in,right=0.6in,bottom=0.5in]{geometry} % Document margins
\usepackage[fixed]{fontawesome5}
\usepackage[breaklinks=true,colorlinks,bookmarks=false]{hyperref}
\hypersetup{
	colorlinks=true,
	linkcolor=gray,
	filecolor=blue,
	citecolor=blue,
	urlcolor=black,
}

\usepackage{enumitem}
\setlist[itemize]{noitemsep, topsep=0pt}

% Publications grouped by type (lists generated by scripts/build_cv.py in publication_groups.tex)
\usepackage[resetlabels]{multibib}
\newcites{jour,conf,thes,pre}{{\vskip -1.1cm},{\vskip -1.1cm},{\vskip -1.1cm},{\vskip -1.1cm}}
\makeatletter
\newcommand\pubprefix[1]{\renewcommand\@biblabel[1]{[#1##1]}}
\makeatother
\newcommand\pubgroup[1]{\par\needlines{7\baselineskip}\vspace{0.2em}\noindent{\small\textbf{#1}}\par}

% Discrete icon-based tags for Certificate/Results/GitHub links
% Allow slightly looser lines instead of text running past the right margin
\setlength{\emergencystretch}{3em}

% \nolinebreak keeps each icon on the same line as the preceding word (never alone on a new line)
\newcommand{\certlink}[1]{\nolinebreak\,\href{#1}{\footnotesize\faCertificate}}
\newcommand{\resultlink}[1]{\nolinebreak\,\href{#1}{\footnotesize\faChartBar}}
\newcommand{\ghlink}[1]{\nolinebreak\,\href{#1}{\footnotesize\faGithub}}
\newcommand{\weblink}[1]{\nolinebreak\,\href{#1}{\footnotesize\faGlobe}}

\newcommand*{\boldname}[3]{%
  \def\lastname{#1}%
  \def\firstname{#2}%
  \def\firstinit{#3}}
\boldname{}{}{}

%---------------------------------------------------------------------------------------
\name{Joao P. A. Dantas} % Your name

% resume.cls prints the 3rd \address first, then the 1st and the 2nd
\address{ \href{https://www.linkedin.com/in/jpdantas/}{\faLinkedin linkedin.com/in/jpdantas} \href{https://www.joaopadantas.com/}{\faGlobeAmericas joaopadantas.com}}

\address{ \href{https://scholar.google.com/citations?user=RjeTT9cAAAAJ}{\faGraduationCap Google Scholar} \href{https://orcid.org/0000-0003-0300-8027}{\faOrcid orcid.org/0000-0003-0300-8027}}

\address{\href{mailto:jpdantas@ita.br}{\faEnvelope jpdantas@ita.br} \href{https://wa.me/5512991335303}{\faPhoneSquare* +55 12 99133 5303}}
% \href{https://github.com/jpadantas}{\faGithub github.com/jpadantas}
\begin{document}
\small

%----------------------------------------------------------------------------------------
%	EDUCATION SECTION
%----------------------------------------------------------------------------------------

\begin{rSection}{Education}

{\bf \href{http://www.ita.br/}{Aeronautics Institute of Technology}}, São José dos Campos, Brazil
\hfill {\em Jan 2019 -- Jun 2025} \\
Ph.D. in Electronic and Computer Engineering \\
Cum. GPA: 9.5 / 10.0

{\bf \href{https://www.cmu.edu/}{Carnegie Mellon University}}, Pittsburgh, PA
\hfill {\em Nov 2021 -- Jan 2023} \\
Ph.D. Exchange Program in Robotics

{\bf \href{http://www.ita.br/}{Aeronautics Institute of Technology}}, São José dos Campos, Brazil
\hfill {\em Aug 2016 -- Jan 2019} \\
M.Sc. in Electronic and Computer Engineering \\
Cum. GPA: 8.5 / 10.0

{\bf \href{https://www.stonybrook.edu/}{Stony Brook University}}, Stony Brook, NY
\hfill {\em Jul 2014 -- Jul 2015} \\
One-year Undergraduate Exchange Program in Mechanical Engineering\\
Cum. GPA: 3.93 / 4.00

{\bf \href{http://www.ita.br/}{Aeronautics Institute of Technology}}, São José dos Campos, Brazil
\hfill {\em Jan 2011 -- Dec 2015} \\
B.Sc. in Mechanical-Aeronautical Engineering\\
Cum. GPA: 8.7 / 10.0
\end{rSection}

%----------------------------------------------------------------------------------------
%	RESEARCH EXPERIENCE
%----------------------------------------------------------------------------------------

\begin{rSection}{Research Experience}

\begin{rSubsection}{SETMV --- Virtual Military Shooting Range Simulator}{Jul 2026 -- Present}{Department of Aerospace Science and Technology (DCTA)}{São José dos Campos, Brazil}
\item Virtual reality and 3D simulation for military shooting instruction, with performance analytics comparing simulated training with live-fire benchmarks to reduce training costs and ammunition usage.
\end{rSubsection}

\begin{rSubsection}{C2 --- Command and Control through Aerospace Data Fusion}{Mar 2025 -- Present}{Institute for Advanced Studies (IEAv)}{São José dos Campos, Brazil}
\item Advanced Command and Control (C2) frameworks leveraging massive aerospace data fusion and high-performance computing to support operational planning in complex defense scenarios.
\end{rSubsection}

\begin{rSubsection}{\href{http://www.asa.dcta.mil.br/}{ASA --- Aerospace Simulation Environment}}{Dec 2020 -- Present}{Institute for Advanced Studies (IEAv)}{São José dos Campos, Brazil}
\item Evaluation of military operational scenarios and tactical protocols for the Brazilian Air Force using the custom-built, object-oriented \textbf{ASA} simulation framework.
\end{rSubsection}

\begin{rSubsection}{Co-development of a Bespoke Air-to-Air Weapon Simulation Environment}{Mar 2024 -- Dec 2025}{Institute for Advanced Studies (IEAv)}{São José dos Campos, Brazil}
\item Offset-agreement cooperation between IEAv and MBDA to co-develop an air-to-air weapon simulation environment integrated with IEAv simulation capabilities and to train the Brazilian team.
\end{rSubsection}

\begin{rSubsection}{Safe and Seamless Operation of Manned and Unmanned Aircraft}{Nov 2021 -- Jan 2023}{\href{https://www.ri.cmu.edu/}{Robotics Institute, Carnegie Mellon University}}{Pittsburgh, PA}
\item Worked with \textbf{Prof.~Dr.~Sebastian Scherer} and \textbf{Prof.~Dr.~Jean Oh} on AI for autonomous aircraft integration into manned traffic, focusing on vision-based detection and trajectory prediction.
\end{rSubsection}

\begin{rSubsection}{ITACOPTER 1 --- Hexacopter for Civil Construction Monitoring}{Jul 2015 -- Dec 2015}{\href{http://www.ita.br/}{Aeronautics Institute of Technology (ITA)}}{São José dos Campos, Brazil}
\item Worked with \href{https://scholar.google.com.br/citations?user=K7tHSrYAAAAJ&hl=BR}{\textbf{Prof.~Dr.~David~Antonio~dos~Santos}} on the design and flight testing of an unmanned hexacopter for infrastructure monitoring.
\end{rSubsection}

\begin{rSubsection}{Stochastic Transport Processes at Interfaces}{May 2015 -- Aug 2015}{\href{https://www.stonybrook.edu/}{Stony Brook University}}{Stony Brook, NY}
\item Worked with {\href{https://me.stonybrook.edu/people/faculty/Colosqui_Carlos.html}{\textbf{Prof.~Dr.~Carlos~Colosqui}}} on the numerical modeling of Brownian Motion and colloidal particle dynamics at micro- and nanoscale interfaces.
\end{rSubsection}

\begin{rSubsection}{Automatic Data Acquisition in Plasma Tunnel Tests}{Aug 2012 -- Jul 2013}{\href{http://www.ita.br/}{Aeronautics Institute of Technology (ITA)}}{São José dos Campos, Brazil}
\item Worked with \href{https://scholar.google.com.br/citations?user=gu-xROMAAAAJ&hl=pt-BR}{\textbf{Prof.~Dr.~Gilberto~Petraconi}} on development of automatic data acquisition systems for plasma environments and evaluation of thermo-structural composite materials for aerospace protective coatings.
\end{rSubsection}

\begin{rSubsection}{Thermionic Cathode Discharge in Low Pressure}{Aug 2011 -- Jul 2012}{\href{http://www.ita.br/}{Aeronautics Institute of Technology (ITA)}}{São José dos Campos, Brazil}
\item Worked with \href{https://scholar.google.com.br/citations?user=gu-xROMAAAAJ&hl=pt-BR}{\textbf{Prof.~Dr.~Gilberto~Petraconi}} on the experimental analysis of thermionic cathode discharge and plasma physics in low-pressure environments for ion source efficiency improvement.
\end{rSubsection}

% \begin{rSubsection}{Plasma Physics and Aerospace Materials}{Aug 2011 -- Jul 2013}{\href{http://www.ita.br/}{Aeronautics Institute of Technology (ITA)}}{São José dos Campos, Brazil}
% \item Worked with \href{https://scholar.google.com.br/citations?user=gu-xROMAAAAJ&hl=pt-BR}{\textbf{Prof.~Dr.~Gilberto~Petraconi}} on the experimental analysis of thermionic cathode discharge for ion source improvement, and the development of automatic data acquisition systems for plasma environments to characterize thermo-structural composite materials used in aerospace protective coatings.
% \end{rSubsection}

\end{rSection}

%----------------------------------------------------------------------------------------
%	PROFESSIONAL EXPERIENCE
%----------------------------------------------------------------------------------------

\begin{rSection}{Professional Experience}

\begin{rSubsection}{\href{https://ieav.dcta.mil.br/}{Institute for Advanced Studies (IEAv)}}{Jan 2016 -- Present}{Research Engineer}{São José dos Campos, Brazil}
    \item Research on military \textbf{Modeling \& Simulation}, \textbf{Machine Learning}, and data science for decision support systems, building simulation environments and analytical tools for complex aerospace scenarios.
\end{rSubsection}

\begin{rSubsection}{\href{https://www.fab.mil.br/index.php}{Brazilian Air Force (FAB)}}{Feb 2013 -- Present}{Captain Engineer Officer}{São José dos Campos, Brazil}
    \item \textbf{Captain} (Dec 2022 -- Present): AI researcher, coordinator of defense offset programs, and firearms instructor.
    \item \textbf{First Lieutenant} (Dec 2015 -- Dec 2022): data science research on defense scenario analysis and operational modeling.
    \item \textbf{Officer Candidate} (Feb 2013 -- Dec 2015): ITA research on aerial robotics, plasma physics, and advanced materials.
\end{rSubsection}

\begin{rSubsection}{\href{https://site.coc.com.br/}{COC Educational System}}{Feb 2017 -- Jul 2021}{Mathematics Teacher and Scientific Olympiads Coordinator}{São José dos Campos, Brazil}
    \item Orchestrated preparatory classes for Scientific Olympiads, guiding students to secure \textbf{158 medals} in 46 national and international competitions.
\end{rSubsection}

\begin{rSubsection}{\href{https://emergebrasil.in/en/}{Emerge Brazil}}{Oct 2016 -- Dec 2017}{Co-founder and Head of Communications}{São Paulo, Brazil}
    \item Bridged scientific research with industry, mapping over \textbf{300 technologies} and facilitating \textbf{60} technology-industry partnerships in the Brazilian market.
\end{rSubsection}

\end{rSection}


%----------------------------------------------------------------------------------------
%	SKILLS SECTION
%----------------------------------------------------------------------------------------

\begin{rSection}{Skills}
\footnotesize
\begin{tabular}{ @{} >{\bfseries}l @{\hspace{6ex}} l }
Natural Languages & Native Portuguese $\diamond$ Advanced English
\\
Programming & Python $\diamond$ C++ $\diamond$ R $\diamond$ SQL $\diamond$ MATLAB $\diamond$ \LaTeX
\\
DL Frameworks & TensorFlow $\diamond$ Keras $\diamond$ PyTorch
\\
Toolkits & NumPy $\diamond$ Pandas $\diamond$ Matplotlib $\diamond$ Scikit-Learn $\diamond$ SciPy $\diamond$ Git $\diamond$ Docker $\diamond$ Linux
\end{tabular}

\end{rSection}


%----------------------------------------------------------------------------------------
%	PUBLICATIONS
%----------------------------------------------------------------------------------------

\begin{rSection}{Publications}
\footnotesize
% Generated by scripts/build_cv.py from publications.json - do not edit by hand.
\pubgroup{Journal Articles}
\pubprefix{J}
\nocitejour{nunes2026stopping,dantas2026poker,silva2025asafg,dantas2025deep,dantas2025autonomous,costa2025simulation,dantas2025asasimaas,murca2025open,patrikar2025image,navarro2024sorts,viscardi2024advanced,dantas2022machine,costa2022formation,dantas2020modelo,dantas2018analise,dantas2017teoria}
\bibliographystylejour{unsrt}
\bibliographyjour{publications}

\pubgroup{Conference Papers}
\pubprefix{C}
\nociteconf{nunes2026haltsim,dantas2026offline,dantas2026harbor,dantas2025offline,araujo2025analise,lima2025algoritmo,lima2025droneswarm2d,dantas2025asa,dantas2024asapy,dantas2024autonomous,dantas2024loyal,lima2024aprendizado,dantas2023real,gobi2023ambiente,dantas2023autonomous,dantas2022supervised,patrikar2022challenges,dantas2021weapon,dantas2021engagement,dantas2021enhanced,dantas2013aquisicao,dantas2012estudo}
\bibliographystyleconf{unsrt}
\bibliographyconf{publications}

\pubgroup{Theses}
\pubprefix{T}
\nocitethes{dantas2025simulation,dantas2018apoio,dantas2015hexacoptero}
\bibliographystylethes{unsrt}
\bibliographythes{publications}

\pubgroup{Preprints}
\pubprefix{P}
\nocitepre{costa2023analysis}
\bibliographystylepre{unsrt}
\bibliographypre{publications}

\end{rSection}

%----------------------------------------------------------------------------------------
%	SOFTWARE
%----------------------------------------------------------------------------------------

\begin{rSection}{Software}
\footnotesize
\textbf{Dec 2025:} DroneSwarm2D --- Autonomous drone swarm simulator for distributed defensive tactics. (INPI, 2025).
\ghlink{https://github.com/lucasll37/DroneSwarm2D-lib}\\
\textbf{Dec 2025:} AsaGym --- Reinforcement Learning library for aerospace scenarios evaluation. \mbox{(INPI BR 51 2026 003324-1).}
\certlink{https://www.joaopadantas.com/certificates/2026_AsaGym_Software_Registry.pdf}
\ghlink{https://github.com/ASA-Simulation/asa-gym}\\
\textbf{Dec 2025:} AsaFG --- Human-in-the-loop integration module for air combat simulations. \mbox{(INPI BR 51 2026 003322-5).}
\certlink{https://www.joaopadantas.com/certificates/2026_AsaFG_Software_Registry.pdf}
\ghlink{https://github.com/ASA-Simulation/AsaFG}\\
\textbf{Oct 2024:} \href{https://www.asa.dcta.mil.br/}{ASA (Ambiente de Simulação Aeroespacial)} --- Simulation framework for modeling military scenarios. \mbox{(INPI BR 51 2024 002174-4).}
\certlink{https://www.joaopadantas.com/certificates/2024_ASA_Software_Registry.pdf}
\weblink{https://www.asa.dcta.mil.br/}\\
\textbf{Jun 2023:} AI Pilot --- Close-proximity safe and seamless operation of manned and unmanned aircraft in shared airspace. \mbox{(CMU Tech Transfer 2023-102).}
\certlink{https://www.joaopadantas.com/certificates/2023_CMU_AIPilot_Software_Registry.pdf}\\
\textbf{2025:} \href{https://github.com/castacks/TartanAviation.git}{TartanAviation} --- Large-scale terminal airspace operations datasets.
\ghlink{https://github.com/castacks/TartanAviation.git}
\weblink{https://theairlab.org/aitf/}\\
\textbf{2024:} \href{https://github.com/ASA-Simulation/AsaPy}{AsaPy} --- A Python library for aerospace simulation analysis and visualization.
\ghlink{https://github.com/ASA-Simulation/AsaPy}
\end{rSection}

%----------------------------------------------------------------------------------------
%	SUPERVISION (ORIENTATIONS)
%----------------------------------------------------------------------------------------

\begin{rSection}{Supervision}
\footnotesize
\begin{enumerate}
    \item \textbf{Artur Ligieri Nunes}. Stop Criteria in Defense Simulations: Comparative Analysis and Statistical Approach Proposal. 2025. M.Sc. Thesis (Electronic and Computer Engineering) --- Aeronautics Institute of Technology. (\textit{Ongoing}).
    \item \textbf{Lucas Silva Lima}. Proposta de Algoritmo Distribuído para Enxame de Drones Defensivos de Pequeno Porte Baseado em Inteligência Artificial. 2025. M.Sc. Thesis (Electronic and Computer Engineering) --- Aeronautics Institute of Technology.
    \item \textbf{Lucas Silva Lima}. Proposta de Algoritmo Distribuído para Enxame de Drones Defensivos de Pequeno Porte Baseado em Planejamento. 2025. B.Sc. Project (Computer Engineering) --- Aeronautics Institute of Technology.
    \item \textbf{Victor Luiz Dorea Beirão Costa}. Apoio à Decisão no Planejamento Baseado em Capacidades: Uma Análise de Cenários de Defesa no Ambiente de Simulação Aeroespacial. 2024. B.Sc. Project (Mechanical-Aeronautical Engineering) --- Aeronautics Institute of Technology.
    \item \textbf{Gabriel Henrique Gobi} and \textbf{Pedro Lustosa Rege Botelho}. AsaGym: A Python Library for Reinforcement Learning in Aerospace Operational Scenarios. 2023. B.Sc. Project (Computer and Electronic Engineering) --- Aeronautics Institute of Technology.
\end{enumerate}
\end{rSection}

%----------------------------------------------------------------------------------------
%	HONORS & RECOGNITIONS and FELLOWSHIPS & GRANTS (generated from awards.json by scripts/build_cv.py)
%----------------------------------------------------------------------------------------

% Generated by scripts/build_cv.py from awards.json - do not edit by hand.
\begin{rSection}{Honors \& Recognitions}
\footnotesize\setlength{\parskip}{0pt}\interlinepenalty=10000
\textbf{Dec 2025:} 3rd place in the Simulation Challenge at the \href{https://meetings.informs.org/wordpress/wsc2025/}{IEEE Winter Simulation Conference 2025} (\$1,000).
\certlink{https://www.joaopadantas.com/certificates/2025_WSC_Simulation_Challenge.pdf}
\resultlink{https://meetings.informs.org/wordpress/wsc2025/simulation-challenge/}

\textbf{Nov 2025:} 1st place in the \href{https://www.gov.br/defesa/pt-br}{Brazilian Ministry of Defense Doctoral Award} 2025 for the best Ph.D. thesis in defense science and technology (R\$4,000).
\certlink{https://www.joaopadantas.com/certificates/2025_MD_PhD_Award.pdf}

\textbf{Dec 2024:} WSC Diversity Award at the \href{https://meetings.informs.org/wordpress/wsc2024/}{IEEE Winter Simulation Conference 2024} in Orlando, USA.
\certlink{https://www.joaopadantas.com/certificates/2024_WSC_Diversity_Award.pdf}
\resultlink{https://connect.informs.org/simulation/awards/wsc-diversity-award/awardees}

\textbf{Dec 2024:} Runner-up Best Poster Award at the \href{https://www.lacoro.org/2024/}{LACORO Summer School on Robotics 2024} in Chile.
\certlink{https://www.joaopadantas.com/certificates/2024_LACORO_Best_Poster.pdf}
\resultlink{https://www.lacoro.org/2024/awardees/}

\textbf{Nov 2024:} 2nd place in the \href{https://prc-data-challenge.netlify.app/}{Performance Review Commission (PRC) Data Challenge 2024} with an open \href{https://github.com/PRC-Data-Challenge-2024/team_tiny_rainbow}{Machine Learning model to infer flight Actual TakeOff Weight (ATOW)}.
\resultlink{https://prc-data-challenge.netlify.app/\#and-the-winners-are}

\textbf{Jun 2024:} ML and Systems Rising Star 2024, awarded by \href{https://mlcommons.org/2024/06/2024-mlc-rising-stars/}{MLCommons}.
\resultlink{https://mlcommons.org/2024/06/2024-mlc-rising-stars/}

\textbf{Oct 2023:} 3rd place in the \href{https://comp.ita.br/dsc/edicoes/2023/}{Data Science Challenge 2023} for predicting the Estimated Landing Time (ELDT) of flights.
\resultlink{https://comp.ita.br/dsc/edicoes/2023/}

\textbf{Jun 2023:} Best PhD Colloquium Award at the \href{https://sigsim.acm.org/conf/pads/2023/}{Conference on Principles of Advanced Discrete Simulation (PADS) 2023}.
\certlink{https://www.joaopadantas.com/certificates/2023_PADS_Best_PhD_Colloquium.pdf}

\textbf{Oct 2021:} 3rd place in the \href{https://www.kaggle.com/c/kddbr-2021}{5th Brazilian Competition on Knowledge Discovery in Databases (KDD-BR)}, predicting solutions of the Travelling Salesman Problem (R\$200).
\resultlink{https://www.kaggle.com/competitions/kddbr-2021/leaderboard}

\textbf{Oct 2021:} Military Bronze Medal, awarded by the \href{https://www.fab.mil.br/index.php}{Brazilian Air Force}.
\certlink{https://www.joaopadantas.com/certificates/2021_Military_Bronze_Medal.pdf}

\textbf{Dec 2019:} Ambassador Award of Excellence 2019, \href{https://iymc.info/en/jodantas}{International Youth Math Challenge (IYMC)}.
\certlink{https://www.joaopadantas.com/certificates/2019_IYMC_Ambassador.pdf}

\textbf{Oct 2019:} 6th place (39 competitors) in the \href{https://www.kaggle.com/c/kddbr-2019}{3rd Brazilian Competition on Knowledge Discovery in Databases (KDD-BR)}, predicting flight delays for a Brazilian airline.
\resultlink{https://www.kaggle.com/competitions/kddbr-2019/leaderboard}

\textbf{May 2019:} 4th place in the \href{https://www.kaggle.com/competitions/data-science-challenge-at-eef-2019/}{Data Science Challenge at Engineering Education for the Future (EEF) 2019}, predicting real estate prices.
\resultlink{https://www.kaggle.com/competitions/data-science-challenge-at-eef-2019/leaderboard}

\textbf{Dec 2015:} Honorable Mention in the Department of Humanities, \href{https://www.ita.br/}{Aeronautics Institute of Technology}.
\certlink{https://www.joaopadantas.com/certificates/2015_ITA_Humanities_Honorable_Mention.pdf}

\textbf{Mar 2014:} Military Merit Award, \href{https://www.fab.mil.br/index.php}{Brazilian Air Force} (top 25\% behavioral grades).
\certlink{https://www.joaopadantas.com/certificates/2014_Military_Merit_Award.pdf}

\textbf{Nov 2011:} 1st place overall in the \href{https://cporaersj.dcta.mil.br/}{Brazilian Air Force Reserve Officer Training Corps (CPORAER-SJ)} (120 candidates).
\certlink{https://www.joaopadantas.com/certificates/2011_CPORAER_SJ_1st_Place.pdf}

\textbf{Dec 2010:} ITA Entrance Exam Approval at the \href{https://www.ita.br/}{Aeronautics Institute of Technology} (120 admitted out of 7,627 applicants), one of Brazil's most competitive entrance exams.
\certlink{https://www.joaopadantas.com/certificates/2010_ITA_Entrance_Exam.pdf}

\textbf{Jan 2009:} Academic Merit Award, \href{https://www.fab.mil.br/index.php}{Brazilian Air Force} (top 25\% military grades).

\textbf{Oct 2007:} Bronze Medal in the \href{https://www.oba.org.br/}{10th Brazilian Olympiad of Astronomy and Astronautics (OBA)} (187,726 competitors).
\certlink{https://www.joaopadantas.com/certificates/2007_OBA_Bronze_Medal.pdf}

\textbf{Dec 2006:} Honorable Mention in the XIV Ceara Science Olympiad --- Biology (OCC), \href{https://www.ufc.br/pt}{Federal University of Ceara (UFC)}.
\certlink{https://www.joaopadantas.com/certificates/2006_UFC_Science_Biology_Olympiad.pdf}

\textbf{Dec 2006:} Gold Medal in the Internal Science Olympiads 2006 (Math, Physics, Biology), \href{https://www.fariasbrito.com.br/}{Farias Brito High School}.
\certlink{https://www.joaopadantas.com/certificates/2006_Farias_Brito_Internal_Olympiads.pdf}

\textbf{Nov 2003:} Honor to Merit (three 2nd places, one 5th place) in the Internal Mathematics Olympiad, \href{https://www.christus.com.br/}{Christus Middle School}.
\certlink{https://www.joaopadantas.com/certificates/2003_Christus_Internal_Math_Olympiad.pdf}

\textbf{Oct 2003:} 2nd place in the 3rd Christus Exhibition of Science and Technology (EXCETEC), category ``Exact, Social and Commercial Sciences''.
\certlink{https://www.joaopadantas.com/certificates/2003_EXCETEC_2nd_Place.pdf}
\end{rSection}

\begin{rSection}{Fellowships \& Grants}
\footnotesize\setlength{\parskip}{0pt}\interlinepenalty=10000
\textbf{May 2026:} DAAD AInet Fellow (AI in Neuroscience), \href{https://www.daad.de/en/the-daad/postdocnet/}{German Academic Exchange Service (DAAD)}, Postdoctoral Networking Tour in Artificial Intelligence (Postdoc-NeT-AI).
\resultlink{https://www.daad.de/en/the-daad/postdocnet/fellows/fellows/\#Dantas}

\textbf{Dec 2025:} Travel Grant (\$1,000) to attend the \href{https://meetings.informs.org/wordpress/wsc2025/}{IEEE Winter Simulation Conference 2025} in Seattle, USA.

\textbf{Dec 2025:} Travel Grant (\$1,500), \href{https://www.latinxinai.org/}{LatinX in AI Coalition}, for the \href{https://www.latinxinai.org/neurips-2025}{LatinX in AI Research Workshop} at \href{https://neurips.cc/Conferences/2025}{NeurIPS 2025}, San Diego, USA.

\textbf{Jun 2025:} DAC Young Fellow for the \href{https://dac.com/}{Design Automation Conference (DAC) 2025}.
\certlink{https://www.joaopadantas.com/certificates/2025_DAC_Young_Fellow.pdf}

\textbf{Jun 2025:} Travel Grant (€2,900), \href{https://cisb.org.br/}{CISB} Call 15/2025, for the \href{https://ft2025.se/}{2025 Swedish Aerospace Congress}.
\resultlink{https://cisb.org.br/wp-content/uploads/2025/06/Disclosure-of-results-Call-CISB-15-2024.pdf}

\textbf{Dec 2024:} Travel Grant (\$1,000) to attend the \href{https://meetings.informs.org/wordpress/wsc2024/}{IEEE Winter Simulation Conference 2024} in Orlando, USA.

\textbf{Dec 2024:} Travel Grant (\$500) to attend the \href{https://www.lacoro.org/2024/}{LACORO Summer School on Robotics 2024} in Chile.

\textbf{Jun 2024:} Travel Grant (\$1,000) to attend the \href{https://sigsim.acm.org/conf/pads/2024/}{2024 Conference on Principles of Advanced Discrete Simulation (PADS)}.

\textbf{May 2024:} Travel Grant (\$3,000) to attend the \href{https://disarmament.unoda.org/en}{UNODA-SIPRI Workshop on Responsible AI for Peace and Security} in Brussels.
\certlink{https://www.joaopadantas.com/certificates/2024_UNODA_SIPRI_Award.pdf}

\textbf{May 2024:} Travel Grant (€3,150), \href{https://cisb.org.br/}{CISB} Call 13/2024, for the \href{https://web.archive.org/web/20240419125501/https://portal.waraps.org/page/wara-ps-demonstration-week}{WARA‑PS Demonstration Week 2024}.
\resultlink{https://cisb.org.br/wp-content/uploads/2024/05/Call-CISB-132024-Disclosure-of-results.pdf}

\textbf{Jun 2023:} Travel Grant (€3,200), \href{https://cisb.org.br/}{CISB} Call 10/2022, for an AI researchers delegation in Sweden.
\resultlink{https://cisb.org.br/images/pdf/Disclosure-of-results-CISB-Call-102022.pdf}

\textbf{Jun 2023:} Travel Grant (\$1,000) to attend the \href{https://sigsim.acm.org/conf/pads/2023/}{Conference on Principles of Advanced Discrete Simulation (PADS) 2023} in Orlando, USA.
\certlink{https://www.joaopadantas.com/certificates/2023_PADS_Travel_Award.pdf}

\textbf{Apr 2023:} Travel Grant (\$2,500) to attend \href{https://brazilatsiliconvalley.com/}{Brazil at Silicon Valley 2023} in Santa Clara, USA.

\textbf{Jun 2022:} Merit Scholarship, \href{https://www.estudar.org.br/comunidade-de-lideres/}{Fundação Estudar}, awarded for outstanding trajectory and academic potential (30 of 33,876 applicants).
\resultlink{https://issuu.com/estudar/docs/fe_relatorio_2022_digital/20}

\textbf{Oct 2016:} \href{https://www.programaprolider.com.br/}{Prolider} Program Fellow 2016, selected by \href{https://www.institutofour.org/}{Instituto Four} (29 of 2,013 applicants).
\certlink{https://www.joaopadantas.com/certificates/2016_Prolider_Fellow.pdf}

\textbf{Aug 2014 -- May 2015:} Full Scholarship ($\approx$\$75,000) from Brazil's Ministry of Education for a one-year study program at Stony Brook University, College of Engineering.
\certlink{https://www.joaopadantas.com/certificates/2014_Stony_Brook_Scholarship.pdf}

\textbf{Aug 2012 -- Jul 2013:} Undergraduate Research Fellowship (\href{http://memoria.cnpq.br/pibic}{PIBIC}), \href{https://www.gov.br/cnpq/pt-br}{CNPq}, at the Aeronautics Institute of Technology, for the project ``Automatic Data Acquisition in Plasma Tunnel Tests''.
\certlink{https://www.joaopadantas.com/certificates/2013_PIBIC_Fellowship_Aquisicao.pdf}

\textbf{Aug 2011 -- Jul 2012:} Undergraduate Research Fellowship (\href{http://memoria.cnpq.br/pibic}{PIBIC}), \href{https://www.gov.br/cnpq/pt-br}{CNPq}, at the Aeronautics Institute of Technology, for the project ``Thermionic Cathode Discharge in Low Pressure''.
\certlink{https://www.joaopadantas.com/certificates/2012_PIBIC_Fellowship_Estudo.pdf}

\textbf{Jan 2004 -- Dec 2006:} Full Scholarship ($\approx$\$30,000) for three years at \href{https://www.fariasbrito.com.br/}{Farias Brito High School}.
\end{rSection}


%----------------------------------------------------------------------------------------
%	VOLUNTEERING
%----------------------------------------------------------------------------------------

\begin{rSection}{Volunteering}

\begin{rSubsection}{\href{https://www.estudar.org.br/}{Fundação Estudar}}{Jan 2025 -- Dec 2025}{Ambassador}{São Paulo, Brazil}
\item Engaging with a network of high-achieving young leaders, supporting initiatives to promote education and career development for Brazilian students.
\end{rSubsection}

\begin{rSubsection}{\href{https://iaac.space/}{International Astronomy and Astrophysics Competition}}{Apr 2020 -- Jul 2021}{Ambassador}{São José dos Campos, Brazil}
\item Inspiring youths for astronomy and astrophysics and encouraging participation in the competition.
\end{rSubsection}

\begin{rSubsection}{\href{https://iymc.info/}{International Youth Math Challenge}}{Jul 2019 -- Jul 2021}{Ambassador}{São José dos Campos, Brazil}
\item Informing schools and encouraging students and youths to participate in the competition.
\end{rSubsection}

\begin{rSubsection}{\href{https://alphalumen.org.br/}{Alpha Lumen Institute}}{Jun 2012 -- May 2013}{Mathematics Teacher}{São José dos Campos, Brazil}
\item Teacher of mathematics in the preparatory class for the ITA entrance exam.
\end{rSubsection}

\end{rSection}

%----------------------------------------------------------------------------------------
%	OPERATIONAL COURSES
%----------------------------------------------------------------------------------------

\begin{rSection}{Operational Courses}
\begin{itemize}
    \item \textbf{Emergency Parachute Jump Training (ISE)} -- Brazilian Air Force (2008).
    \item \textbf{Firearms Instructor Course} -- Brazilian Air Force (2023).
    \item \textbf{Open Water Diver Course (Certification No. 26030I6411)} -- Professional Association of Diving Instructors (PADI) (2026).
    \item \textbf{Military Remotely Piloted Aircraft Operator Course} -- Brazilian Air Force (2026).
\end{itemize}
\end{rSection}

%----------------------------------------------------------------------------------------
%	ADDITIONAL ACTIVITIES
%----------------------------------------------------------------------------------------

\begin{rSection}{Additional Activities}
\begin{itemize}
\item \textbf{Sports:} Basketball player at Air Force Preparatory School and Air Force Academy (2007--2009). Amateur bodybuilder (2014--2015) -- \href{https://contests.npcnewsonline.com/contests/2015/npc_northeast_summer_classic/2cdz/joao_dantas}{The National Physique Committee}. Active Brazilian Jiu-Jitsu practitioner (blue belt, IBJJF).
\item \textbf{Outdoors and Interests:} Orienteering, trekking, reading (science fiction), chess, and traveling.
\end{itemize}
\end{rSection}

\today

\end{document}
